% =======================================================================================
\section{Network Architecture}
\label{app:architecture}
% =======================================================================================

The velocity field is a transformer with two token streams, one for the sources and one for
the data, that attend to each other in every block: the multimodal diffusion transformer
(MMDiT) of \citet{esser2024scaling}, with the adaptive-normalisation conditioning of
\citet{peebles2023scalable}. \Cref{fig:architecture} shows it, and \cref{tab:network} gives
the sizes.

\begin{figure}[h]
  \centering
  \includegraphics[width=\textwidth]{figures/architecture.pdf}
  \caption{The network. The parameter stream (blue) has one token per source and no
  positional embedding, so the network is equivariant to permutations of the sources; the
  data stream (grey) has one token per $2\times2$ patch of the conditioning image. Both
  streams attend jointly in every block, and the flow time and the window frequency modulate
  every normalisation and residual gate. Only the velocity head is used to sample; the two
  auxiliary heads are trained during the first 10\% of the steps.}
  \label{fig:architecture}
\end{figure}

\subsection{Embeddings}

All feed-forward maps are SwiGLU \citep{shazeer2020glu},
$\mathrm{FF}(z) = W_2\bigl(a\odot\mathrm{SiLU}(g)\bigr)$ with $[a;g] = W_1z + b_1$.
\begin{itemize}
  \item \textbf{Sources.} Each source's 11 coordinates are embedded by its own application of
        a shared $\mathrm{FF}: \R^{11}\to\R^D$, giving $S$ tokens with no positional embedding.
  \item \textbf{Data.} The $2\times K\times3$ image is cut into $2\times2$ patches, each
        flattened to 12 numbers and mapped linearly to $\R^D$, giving $2R$ tokens (1 row of
        $K/2$). Sinusoidal embeddings of the row and column index, each passed through an
        $\mathrm{FF}$, are added.
  \item \textbf{Conditioning.} The flow time and the window centre frequency are embedded as
        $c = \mathrm{FF}_t\bigl(E(t)\bigr) + \mathrm{FF}_f\bigl(E(\fw)\bigr)$, with
        $E(\tau) = \mathrm{FF}\bigl([\sin2\pi\omega_k\tau;\ \cos2\pi\omega_k\tau]_{k=1}^{D}\bigr)$
        and $\omega_k = (2\pi)^{-(k-1)/(D-1)}$, angular frequencies from 1 to $2\pi$ per
        unit. For $\fw$ in hertz the phases are small and the embedding is close to linear in
        $\fw$.
\end{itemize}

\subsection{Blocks}

Each of the $N$ blocks updates the source tokens $z^x$ and the data tokens $z^y$. For each
stream $\sigma\in\{x,y\}$, a zero-initialised linear map of $\mathrm{SiLU}(c)$ produces a shift,
a scale and a gate, $(\beta^\sigma,\gamma^\sigma,\zeta^\sigma)$, so that every residual branch
starts as the identity \citep{peebles2023scalable}:
\begin{align}
  h^\sigma &= \mathrm{LN}(z^\sigma)\odot(1+\gamma^\sigma) + \beta^\sigma, \nonumber\\
  [q^\sigma; k^\sigma; v^\sigma] &= W^\sigma_{qkv}h^\sigma,\qquad
  o = \mathrm{Attn}\bigl([q^x; q^y], [k^x; k^y], [v^x; v^y]\bigr), \nonumber\\
  z^\sigma &\leftarrow z^\sigma + \zeta^\sigma\odot W^\sigma_{o}\,o^\sigma, \nonumber\\
  z^\sigma &\leftarrow z^\sigma + \zeta'^\sigma\odot\mathrm{FF}^\sigma\bigl(\mathrm{LN}(z^\sigma)\odot(1+\gamma'^\sigma) + \beta'^\sigma\bigr),
\end{align}
with a second shift, scale and gate $(\beta'^\sigma,\gamma'^\sigma,\zeta'^\sigma)$ for the
feed-forward branch, $\mathrm{LN}$ a layer normalisation without learned affine parameters,
and the feed-forward hidden width $2D$. The attention runs over all $S+2R$ tokens with $H$
heads of dimension 64. Queries and keys are RMS-normalised per head, with no learned scale
\citep{henry2020query,dehghani2023scaling}. The projections are separate for the two streams,
the attention is joint. The blocks are applied with a scan over stacked parameters and
rematerialised in the backward pass \citep{chen2016training}.

\subsection{Heads}

The source tokens pass through a final shift-and-scale modulation and an
$\mathrm{FF}:\R^D\to\R^{2\times11}$, which gives the velocity $\vnet$ and the endpoint
estimate $\hat X_1$ for each source. The data tokens are mapped back to $2\times2$ patches,
which gives the reconstruction $\hat y$.

\subsection{Permutation Equivariance}
\label{app:architecture-equivariance}

Every map applied to a source token is shared across sources, and attention without
positional information is equivariant to permutations of its tokens. Permuting the input
sources therefore permutes the outputs: $\vnet(\sigma X) = \sigma\,\vnet(X)$ for every
$\sigma\in\mathfrak S_S$. With an exchangeable base distribution, the flow of \cref{eq:ode}
maps exchangeable distributions to exchangeable distributions, and the set alignment of
\cref{eq:set-log} gives it a regression target that is itself defined on unordered sets
\citep[compare][]{zaheer2017deep,lee2019set}.

\subsection{Sizes and Precision}

\Cref{tab:network} gives the two widths we trained. The network is the same on every rung,
and only the number of data tokens changes. Parameters are kept in float32, and the forward
and backward passes run in bfloat16 \citep{micikevicius2018mixed,kalamkar2019study}. In
bfloat16, attention uses the fused cuDNN flash-attention kernel \citep{dao2022flashattention}:
on B, materialising the attention logits of a batch, $256\times8$ heads $\times\,2052^2$ in
float32, would need 32\,GiB, which is what made the first B run run out of memory.

\begin{table}[h]
\caption{Network sizes. Parameter counts are exact; tokens are sources plus data patches; step times are for a batch of 256 windows (\cref{app:training-compute}).}
\label{tab:network}
\centering
\small
\begin{tabular}{lrrrrr}
\toprule
width $D$ & heads $H$ & blocks $N$ & parameters & tokens on XS / S / B & A100 time per step, XS / B \\
\midrule
512 & 8 & 8 & 75\,636\,258 & 36 / 132 / 2052 & 38.9\,ms / 1.54\,s \\
768 & 12 & 8 & 170\,077\,474 & 36 / 132 / 2052 & 59.6\,ms / 2.57\,s \\
\bottomrule
\end{tabular}
\end{table}
